\section{Introduction}
\label{sec:introduction}


1.1 Background
•	Increasing relevance of cloud governance in multi-cloud and hybrid-cloud environments 
•	Use of the term landing zone by major cloud providers such as Azure, AWS, GCP, Alibaba Cloud, and IBM Cloud 
•	Importance of structured foundational environments for secure and manageable cloud adoption 
1.2 Problem Statement
•	The concept of a landing zone is not defined uniformly across providers 
•	Differences in terminology, structure, and emphasis make direct comparison difficult 
•	A comparative understanding is needed to clarify the concept in a technical context 
1.3 Aim of the Report
•	To analyze how major cloud providers define and structure landing zones 
•	To identify the main similarities and differences between provider concepts 
•	To formulate a common definition based on the comparison 
1.4 Research Question
•	How do landing zone concepts differ across major cloud providers, and what common definition can be derived from their similarities and differences? 
1.5 Scope and Limitations
•	Focus on conceptual comparison rather than implementation details 
•	Focus on selected providers: Azure, AWS, GCP, Alibaba Cloud, and IBM Cloud 
•	Focus on landing zones as foundational structures for governance, security, and workload enablement 
•	Limitation to publicly available provider documentation and selected technical literature 
1.6 Structure of the Report
•	Chapter 1 introduces the topic, problem, objective, and scope 
•	Chapter 2 reviews the conceptual and technical background 
•	Chapter 3 presents the methodology and comparison framework 
•	Chapter 4 provides the comparative analysis and synthesis 
•	Chapter 5 discusses the significance and limitations of the findings 
•	Chapter 6 concludes the report by answering the research question and summarizing the main outcome 
